\SourceLeaf{034}
\begin{center}
{\large THEORIA\par}
\medskip
{\Large\bfseries COMBINATIONIS OBSERVATIONUM\par}
\medskip
{\large ERRORIBUS MINIMIS OBNOXIAE.\par}
\medskip
{\large\textsc{Pars Posterior.}\par}
\medskip
\rule{.22\linewidth}{.3pt}
\end{center}
\Article{23}
Plures adhuc supersunt disquisitiones, per quas theoria praecedens tum illustrabitur tum ampliabitur.

Ante omnia investigare oportet, num negotium eliminationis, cuius adiumento indeterminatae \(x,y,z\) etc. per \(\xi,\eta,\zeta\) etc. exprimendae sunt, semper sit possibile. Quum multitudo illarum multitudini harum aequalis sit, e theoria eliminationis in aequationibus linearibus constat, illam eliminationem. si \(\xi,\eta,\zeta\) etc. ab invicem independentes sint, certo possibilem fore; sin minus, impossibilem. Supponamus aliquantisper, \(\xi,\eta,\zeta\) etc. non esse ab invicem independentes. sed exstare inter ipsas aequationem identicam
\[
0=F\xi+G\eta+H\zeta+\text{etc.}+K
\]
Habebimus itaque
\[
\begin{aligned}
F\Sigma aa+G\Sigma ab+H\Sigma ac+\text{etc.}&=0\\
F\Sigma ab+G\Sigma bb+H\Sigma bc+\text{etc.}&=0\\
F\Sigma ac+G\Sigma bc+H\Sigma cc+\text{etc.}&=0
\end{aligned}
\]
etc., nec non
\[
F\Sigma al+G\Sigma bl+H\Sigma cl+\text{etc.}=-K
\]
\par\SourceLeaf{035}
Statuendo porro
\[
\left.\begin{aligned}
aF+bG+cH+\text{etc.}&=\theta\\
a'F+b'G+c'H+\text{etc.}&=\theta'\\
a''F+b''G+c''H+\text{etc.}&=\theta''
\end{aligned}\right\}\quad\text{(I)}
\]
etc., eruitur
\[
\begin{aligned}
a\theta+a'\theta'+a''\theta''+\text{etc.}&=0\\
b\theta+b'\theta'+b''\theta''+\text{etc.}&=0\\
c\theta+c'\theta'+c''\theta''+\text{etc.}&=0
\end{aligned}
\]
etc., nec non
\[
l\theta+l'\theta'+l''\theta''+\text{etc.}=-K
\]
Multiplicando itaque aequationes (I) resp. per \(\theta,\theta',\theta''\) etc., et addendo, obtinemus:
\[
0=\theta\theta+\theta'\theta'+\theta''\theta''+\text{etc.}
\]
quae aequatio manifesto consistere nequit, nisi simul fuerit \(\theta=0,\allowbreak\theta'=0,\allowbreak\theta''=0\) etc. Hinc primo colligimus, necessario esse debere \(K=0\). Dein aequationes (I) docent, functiones \(v,v',v''\) etc. ita comparatas esse, ut ipsarum valores non mutentur, si valores quantitatum \(x,y,z\) etc. capiant incrementa vel decrementa ipsis \(F,G,H\) etc. resp. proportionalia, idemque manifesto de functionibus \(V,V',V''\) etc. valebit. Suppositio itaque consistere nequit, nisi in casu tali, ubi vel e valoribus exactis quantitatum \(V,V',V''\) etc. valores incognitarum \(x,y,z\) etc. determinare impossibile fuisset, i. e. ubi problema natura sua fuisset indeterminatum, quem casum a disquisitione nostra exclusimus.
\Article{24}
Denotemus per \(\beta,\beta',\beta''\) etc. multiplicatores, qui eandem relationem habent ad indeterminatam \(y\), quam habent \(\alpha,\alpha',\alpha''\) etc. ad \(x\), puta sit
\[
\begin{aligned}
a[\beta\alpha]+b[\beta\beta]+c[\beta\gamma]+\text{etc.}&=\beta\\
a'[\beta\alpha]+b'[\beta\beta]+c'[\beta\gamma]+\text{etc.}&=\beta'\\
a''[\beta\alpha]+b''[\beta\beta]+c''[\beta\gamma]+\text{etc.}&=\beta''
\end{aligned}
\]
etc., ita ut fiat indefinite
\[
\beta v+\beta'v'+\beta''v''+\text{etc.}=y-B
\]
\par\SourceLeaf{036}
Perinde sint \(\gamma,\gamma',\gamma''\) etc. multiplicatores similes respectu indeterminatae \(z\), puta
\[
\begin{aligned}
a[\gamma\alpha]+b[\gamma\beta]+c[\gamma\gamma]+\text{etc.}&=\gamma\\
a'[\gamma\alpha]+b'[\gamma\beta]+c'[\gamma\gamma]+\text{etc.}&=\gamma'\\
a''[\gamma\alpha]+b''[\gamma\beta]+c''[\gamma\gamma]+\text{etc.}&=\gamma''
\end{aligned}
\]
etc., ita ut fiat indefinite
\[
\gamma v+\gamma'v'+\gamma''v''+\text{etc.}=z-C
\]
et sic porro. Hoc pacto, perinde ut iam in art. 20 inveneramus
\[
\Sigma\alpha a=1,\quad\Sigma\alpha b=0,\quad\Sigma\alpha c=0,\quad\text{etc., nec non}\quad\Sigma\alpha l=-A
\]
etiam habebimus
\[
\begin{aligned}
\Sigma\beta a&=0,\quad\Sigma\beta b=1,\quad\Sigma\beta c=0,\quad\text{etc., atque}\quad\Sigma\beta l=-B\\
\Sigma\gamma a&=0,\quad\Sigma\gamma b=0,\quad\Sigma\gamma c=1,\quad\text{etc., atque}\quad\Sigma\gamma l=-C
\end{aligned}
\]
et sic porro. Nec minus, quemadmodum in art. 20 prodiit \(\Sigma\alpha\alpha=[\alpha\alpha]\), etiam erit
\[
\Sigma\beta\beta=[\beta\beta],\quad\Sigma\gamma\gamma=[\gamma\gamma]\quad\text{etc.}
\]
Multiplicando porro valores ipsorum \(\alpha,\alpha',\alpha''\) etc. (art. 20. IV) resp. per \(\beta,\beta',\beta''\) etc., et addendo, obtinemus
\[
\alpha\beta+\alpha'\beta'+\alpha''\beta''\quad\text{etc.}=[\alpha\beta],\quad\text{sive}\quad\Sigma\alpha\beta=[\alpha\beta]
\]
Multiplicando autem valores ipsorum \(\beta,\beta',\beta''\) etc. resp. per \(\alpha,\alpha',\alpha''\) etc., et addendo, perinde prodit
\[
\alpha\beta+\alpha'\beta'+\alpha''\beta''+\text{etc.}=[\beta\alpha],\quad\text{adeoque}\quad[\alpha\beta]=[\beta\alpha]
\]
Prorsus simili modo eruitur
\[
[\alpha\gamma]=[\gamma\alpha]=\Sigma\alpha\gamma,\quad[\beta\gamma]=[\gamma\beta]=\Sigma\beta\gamma\quad\text{etc.}
\]
\Article{25}
Denotemus porro per \(\lambda,\lambda',\lambda''\) etc. valores functionum \(v,v',v''\) etc., qui prodeunt, dum pro \(x,y,z\) etc. ipsarum valores maxime plausibiles \(A,B,C\) etc. substituuntur, puta
